\section{Introduction}

Features matter.
The last decade of progress on various visual recognition tasks has been based considerably on the use of SIFT \cite{SIFT} and HOG \cite{Dalal05}.
But if we look at performance on the canonical visual recognition task, PASCAL VOC object detection \cite{PASCAL-IJCV}, it is generally acknowledged that progress has been slow during 2010-2012, with small gains obtained by building ensemble systems and employing minor variants of successful methods.

SIFT and HOG are blockwise orientation histograms, a representation we could associate roughly with complex cells in V1, the first cortical area in the primate visual pathway.
But we also know that recognition occurs several stages downstream, which suggests that there might be hierarchical, multi-stage processes for computing features that are even more informative for visual recognition.

\begin{figure}[t!]
\centering
\includegraphics[width=\linewidth]{figures/overview/splash-method.pdf}
\caption{\textbf{Object detection system overview.}
Our system (1) takes an input image, (2) extracts around 2000 bottom-up region proposals, (3) computes features for each proposal using a large convolutional neural network (CNN), and then (4) classifies each region using class-specific linear SVMs.
R-CNN achieves a mean average precision (mAP) of \textbf{53.7\% on PASCAL~VOC~2010}. 
For comparison, \cite{UijlingsIJCV2013} reports 35.1\% mAP using the same region proposals, but with a spatial pyramid and bag-of-visual-words approach.
The popular deformable part models perform at 33.4\%.
On the 200-class \textbf{ILSVRC2013 detection dataset, R-CNN's mAP is 31.4\%}, a large improvement over OverFeat \cite{overfeat}, which had the previous best result at 24.3\%.
}
\figlabel{splash}
\vspace{-1em}
\end{figure}

Fukushima's ``neocognitron'' \cite{fukushima1980neocognitron}, a biologically-inspired hierarchical and shift-invariant model for pattern recognition, was an early attempt at just such a process.
The neocognitron, however, lacked a supervised training algorithm.
Building on Rumelhart \etal \cite{rumelhart86}, LeCun \etal \cite{lecun-89e} showed that stochastic gradient descent via backpropagation was effective for training convolutional neural networks (CNNs), a class of models that extend the neocognitron.

CNNs saw heavy use in the 1990s (\eg, \cite{lecun-98}), but then fell out of fashion, particularly in computer vision, with the rise of support vector machines.
In 2012, Krizhevsky \etal \cite{krizhevsky2012imagenet} rekindled interest
in CNNs by showing substantially higher image classification accuracy on the ImageNet Large Scale Visual Recognition Challenge (ILSVRC) \cite{ILSVRC12,imagenet_cvpr09}.
Their success resulted from training a large CNN on 1.2 million labeled images, together with a few twists on LeCun's CNN (\eg, $\max(x,0)$ rectifying non-linearities and ``dropout'' regularization).

The significance of the ImageNet result was vigorously debated during the ILSVRC 2012 workshop.
%\footnote{For an excellent summary read \url{http://goo.gl/PhTvoE}.}
The central issue can be distilled to the following: To what
extent do the CNN classification results on ImageNet generalize
to object detection results on the PASCAL VOC Challenge?

We answer this question decisively by bridging the gap between image classification and object detection.
This paper is the first to show that a CNN can lead to dramatically higher object detection performance on PASCAL VOC as compared to systems based on simpler HOG-like features.
Achieving this result required solving two problems: localizing objects with a deep network and training a high-capacity model with only a small quantity of annotated detection data.

Unlike image classification, detection requires localizing (likely many) objects within an image.
One approach frames localization as a regression problem.
However, work from Szegedy \etal \cite{szegedy2013deep}, concurrent with our own, 
indicates that this strategy may not fare well in practice (they report a mAP of 30.5\% on VOC 2007 compared to the 58.5\% achieved by our method).
An alternative is to build a sliding-window detector.
CNNs have been used in this way for at least two decades, typically on constrained object categories, such as faces \cite{rowley1998neural,lecun94} and pedestrians \cite{sermanetCVPR13}.
In order to maintain high spatial resolution, these CNNs typically only have two convolutional and pooling layers.
We also considered adopting a sliding-window approach.
However, units high up in our network, which has five convolutional layers, have very large receptive fields ($195 \times 195$ pixels) and strides ($32 \times 32$ pixels) in the input image, which makes precise localization within the sliding-window paradigm an open technical challenge.

Instead, we solve the CNN localization problem by operating within the ``recognition using regions'' paradigm \cite{gu2009recognition}, which has been successful for both object detection \cite{UijlingsIJCV2013} and semantic segmentation \cite{carreira2012cpmc}.
At test time, our method generates around 2000 category-independent region proposals for the input image, extracts a fixed-length feature vector from each proposal using a CNN, and then classifies each region with category-specific linear SVMs.
We use a simple technique (affine image warping) to compute a fixed-size CNN input from each region proposal, regardless of the region's shape.
\figref{splash} presents an overview of our method and highlights some of our results.
Since our system combines region proposals with CNNs, we dub the method R-CNN: Regions with CNN features.

In this updated version of this paper, we provide a head-to-head comparison of R-CNN and the recently proposed OverFeat \cite{overfeat} detection system by running R-CNN on the 200-class ILSVRC2013 detection dataset.
OverFeat uses a sliding-window CNN for detection and until now was the best performing method on ILSVRC2013 detection.
We show that R-CNN significantly outperforms OverFeat, with a mAP of 31.4\% versus 24.3\%.

A second challenge faced in detection is that labeled data is scarce and the amount currently available is insufficient for training a large CNN.
The conventional solution to this problem is to use \emph{unsupervised} pre-training, followed by supervised fine-tuning (\eg, \cite{sermanetCVPR13}).
The second major contribution of this paper is to show that \emph{supervised} pre-training on a large auxiliary dataset (ILSVRC classification), followed by domain-specific fine-tuning on a small dataset (PASCAL), is an effective paradigm for learning high-capacity CNNs when data is scarce.\footnote{In contemporaneously work, \cite{decafICML} shows that the CNN of Krizhevsky \etal, trained on ImageNet, generalizes well to a wide range of datasets and recognition tasks including scene classification, fine-grained sub-categorization, and domain adaptation.}
In our experiments, fine-tuning for detection improves mAP performance by 8 percentage points.
After fine-tuning, our system achieves a mAP of 54\% on VOC 2010 compared to 33\% for the highly-tuned, HOG-based deformable part model (DPM) \cite{lsvm-pami,release5}.

Our system is also quite efficient.
The only class-specific computations are a reasonably small matrix-vector product and greedy non-maximum suppression.
This computational property follows from features that are shared across all categories and that are also two orders of magnitude lower-dimensional than previously used region features (\cf \cite{UijlingsIJCV2013}).

One advantage of HOG-like features is their simplicity: it's easier to understand the information they carry (although \cite{vondrick2013hoggles} shows that our intuition can fail us).
Can we gain insight into the representation learned by the CNN?
Perhaps the densely connected layers, with more than 54 million parameters, are the key?
They are not. 
We ``lobotomized'' the CNN and found that a surprisingly large proportion, 94\%, of its parameters can be removed with only a moderate drop in detection accuracy.
%Perhaps color, which HOG makes very weak use of, is the key?
%Removing color degrades performance only marginally.
Instead, by probing units in the network we  see that the convolutional layers learn a diverse set of rich features (\figref{vislittle}).
%ranging from red-blob detectors to semantically aligned poselet-like \cite{BourdevMalikICCV09} units.

Understanding the failure modes of our approach is also critical for improving it, and so we report results from the detection analysis tool of Hoiem \etal
\cite{hoiem2012diagnosing}.
As an immediate consequence of this analysis, we demonstrate that a simple bounding-box regression method significantly reduces mislocalizations, which are the dominant error mode.

Before developing technical details, we note that because R-CNN operates on regions it is natural to extend it to the task of semantic segmentation.
With minor modifications, we also achieve competitive results on the PASCAL VOC segmentation task, with an average segmentation accuracy of 47.9\% on the VOC 2011 test set.
